\subsection{Experimental Setup}

\begin{table*}[]
\vspace{4pt}
\resizebox{\textwidth}{!}{
\centering
\begin{tabular}{c|c|c|c|c|c|c|c|}
\cline{2-8}
 &
  \multicolumn{6}{c|}{\textbf{Task Time [seconds]}}&\multicolumn{1}{c|}{\textbf{Success Rate}}\\ \hline
\multicolumn{1}{|c|}{\textbf{Planner}} &
  \textbf{S1} &
  \textbf{S2} &
  \textbf{S3} &
  \textbf{S4} &
  \textbf{S5} &
  \textbf{S6} 
  & \textbf{S1 - S6}\\ \hline
\multicolumn{1}{|c|}{PokeRRT*} & 
\textbf{46.43 (21.67)} & 
212.74 (100.01) & 
232.42 (118.63) & 
70.47 (38.82) & 
132.62 (98.27) & 
\textbf{130.01 (84.05)} & 
0.87 (0.31) \\ \hline
\multicolumn{1}{|c|}{PokeRRT} & 
49.69 (21.11) & 
\textbf{196.54 (124.61)} & 
\textbf{167.55 (138.63)} & 
\textbf{64.14 (52.45)} & 
\textbf{116.35 (89.72)}& 
171.77 (103.21) & 
\textbf{0.88 (0.31)} \\ \hline
\multicolumn{1}{|c|}{\begin{tabular}[c]{@{}c@{}}LI-PokeRRT*\end{tabular}} & 
169.87 (71.28) & 
378.90 (122.21) & 
346.70 (111.71) & 
165.91 (44.87) & 
N/A & 
N/A & 
0.53 (0.28) \\ \hline
\multicolumn{1}{|c|}{\begin{tabular}[c]{@{}c@{}}LI-PokeRRT\end{tabular}} & 
123.71 (36.72) & 
328.66 (107.18) & 
322.28 (138.05) &
135.80 (32.91) & 
N/A & 
N/A & 
0.63 (0.18) \\ \hline
\multicolumn{1}{|c|}{\begin{tabular}[c]{@{}c@{}}Push Planner\end{tabular}} & 
122.68 (47.07) & 
284.78 (104.40) & 
249.32 (68.30) & 
117.58 (67.05) & 
N/A & 
N/A & 
0.44 (0.26) \\ \hline\hline
\multicolumn{1}{|c|}{Pick-and-Place} & 
16.29 (3.76) & 
17.71 (3.81) & 
16.14 (3.25) & 
N/A & 
19.15 (4.99) & 
N/A & 
0.67 (0.00)\\ \hline
\end{tabular}
}
\caption{Task times [\textsl{mean (stddev)}] for various planning algorithms in simulation are presented. Success rates are aggregated across all 6 scenarios. Overall, poke planning is faster than push planning and leads to higher success rate than pushing or grasping. Low-impulse (LI) poke planners and the push planner are unsuccessful in S5 and S6 due to action path collisions (S5) or goal region being outside the robot's reachable workspace (S6). The robot is unable to pick or place object in S4 and S6  due to limitations in robot mechanics.\vspace{-12pt}}
\label{tab:eval-sim-time}
\end{table*}

\subsubsection{Scenarios}\label{sec:evaluation-scenarios}

We test multiple algorithms for poking, pushing, and grasping manipulation in six scenarios (\cref{fig:scenarios}).
% \ap{\textsl{PokeRRT} and \textsl{PokeRRT*} expand a robot's reachable workspace, in contrast to pushing and pick-and-place which operate within this region. Pushing also operates under the quasistatic assumption, as a result of which the robot must maintain constant contact with the object. And finally, pick-and-place is also limited by the object dimensions, i.e., if the object is bigger than the gripper width, pick-and-place will fail.}
Our motivating application is the concurrent operation of two robots with potentially non-overlapping reachable regions located in adjacent workcells on a factory floor. Therefore, the objects being manipulated must be accessible by both robots to enable successful collaboration.
Given this motivation, scenarios are designed to test the flexibility of the planner in generating plans for manipulating the object in a planar workspace from a fixed start pose to a goal region, which represents the workspace of a second robot. To ensure that object--obstacle collisions are handled robustly through the replanning strategy, we enforce the following properties in our scenarios: i) the objects used in each scenario are rigid and therefore retain their shape across multiple trials and ii) obstacles are fixed to the table so that collisions do not change the planning configuration space.

Scenarios 1-3 (S1-S3 in \cref{fig:scenarios}) are designed to test baseline manipulation capability through uncluttered (S1) and cluttered (S2, S3) environments. S2 and S3 contain 2 and 4 obstacles, respectively, at fixed poses in the shared robot workspace. The object being manipulated has size $[9, 14, 5]$ cm and mass $m=87$ g. Poke planning, push planning, and pick-and-place will all work in these scenarios since the goal region overlaps with the first robot's reachable workspace.
Scenario 4 (S4 in \cref{fig:scenarios}) contains a bigger cuboid object of size $[11, 18, 11]$ cm and mass $m=112$ g in an obstacle-free workspace. Poke planning and push planning will work in this scenario, however pick-and-place will not since gripper width ($10$ cm) is smaller than the minimum dimension of the cuboid. This scenario represents situations where object properties are incompatible with robot kinematics, i.e., if the object of interest is too heavy or too wide to be grasped or if the end-effector is malfunctional, and the robot needs to formulate an alternate manipulation plan.

Scenario 5 (S5 in \cref{fig:scenarios}) presents two work cells with a tunnel in the center of the table. The first robot cannot push the object to the second robot because the end-effector will collide with the divider along its action path. However, poke planning will succeed due to the short duration of robot--object contact. The first robot is able to grasp the object over the divider in our setup but for larger dividers such as a screen, pick-and-place will fail.
Scenario 6 (S6 in \cref{fig:scenarios}) contains an obstacle-free workspace with non-overlapping reachable regions for each robot. The first robot is able to poke the target object to the second robot's reachable workspace, but cannot push or pick-and-place due to limited reachability. Since pushing operates under the quasitstatic assumption, it requires constant contact between the robot end-effector and the object being manipulated and therefore, manipulation is limited by robot kinematics and reachability. This scenario represents cases where task success is limited by the robot's kinematic characteristics, i.e., if the goal pose is outside the robot's reachable workspace.

\subsubsection{Parameters}

The action vector $\boldsymbol{a_{valid}} = \langle \boldsymbol{p_c}, \boldsymbol{\|\vec{v}_{EE}\|}\rangle$ for our proposed motion planners is generated at runtime. Contour points $p_c$ close to object corners are ignored for clean pokes. The velocity magnitudes $\|\vec{v}_{EE}\|$ are sampled in the $[0.3, 1.0]$ m/s range to increase likelihood of robot joints achieving the commanded operational space velocities. Goal region is defined as the workspace of the second robot, determined as the frequency of inverse kinematics poses achievable on discretized xy-locations on the planar support surface of the object. The empirically-determined replanning threshold is $5cm$ and $10^{\circ}$.

\subsubsection{Algorithms}

\textsl{PokeRRT} and \textsl{PokeRRT*} are compared against several baseline approaches. To evaluate pushing, we use the \textsl{Two-Level Push Planner} presented in \cite{zito2012two} since it also uses simulation in the planning loop to get the next feasible environment state. This work integrates simulation-based forward modeling with sampling-based motion planning to explore the space of feasible pushing actions required to get an object from start to goal. Our baseline approaches, \textsl{Low-Impulse (LI) PokeRRT} and \textsl{Low-Impulse (LI) PokeRRT*}, are designed to show that poking is a more fundamental manipulation skill and encompasses pushing if the applied impulse magnitudes are kept small. They operate similarly to \textsl{PokeRRT} and \textsl{PokeRRT*} but with $\|\vec{v}_{EE}\| = 0.2$ m/s to simulate pushing. Any replanning for a given planner is done using that same planner. Lastly, \textsl{Pick-and-Place} is performed in an open-loop manner with predefined grasps for known objects. An experimental trial fails if the planner does not find a valid plan to the goal region in 240 seconds or if the object falls off the table during execution.\vspace{-4pt}
